\section{Reconstruction limits low-bit decode attention speedups}
\label{sec:analysis}

A low-bit KV cache reduces HBM traffic, but the kernel must still construct floating-point operands fast enough to supply the Tensor Cores.
In existing baselines, reconstruction involves converting integer values into floating-point operands, applying the affine transformation $sq+z$, and optionally performing one additional conversion to match the Tensor Core operand format (e.g., BF16 to FP8).

\subsection{Low-bit baselines fall short of peak HBM throughput}
\label{sec:analysis-gap}

\begin{figure}[t]
  \centering
  \captionsetup[subfloat]{farskip=2pt}
  \includegraphics[width=\columnwidth]{assets/roofline-gap-legend.pdf}
  \par
  \subfloat[H100 (3.35~TB/s peak).\label{fig:roofline-gap-h100}]{%
    \includegraphics{assets/roofline-gap-h100.pdf}}
  \hfill
  \subfloat[B200 (8.00~TB/s peak).\label{fig:roofline-gap-b200}]{%
    \includegraphics{assets/roofline-gap-b200.pdf}}
  \caption{Effective KV access throughput of baseline kernels and \sys{} at $B=128$ and $KV=32\mathrm{K}$ as a percentage of peak HBM throughput.}
  \label{fig:roofline-gap}
\end{figure}


We benchmark FlashInfer~\cite{ye2025flashinfer}, TensorRT-LLM~\cite{nvidia2026trtllm}, FlashAttention-2~\cite{dao2024flashattention2}, FlashAttention-3~\cite{shah2024flashattention3}, FlashAttention-4~\cite{zadouri2026flashattention4}, BitDecoding~\cite{du2026bitdecoding}, KIVI~\cite{liu2024kivi}, Kitty~\cite{xia2025kitty}, SAW-INT4~\cite{jia2026saw}, and OSCAR~\cite{zhou2026oscar} across KV-cache storage formats.
We compute effective KV access throughput by dividing the bytes in one read of the stored values, scales, and offsets by decode time.
\Cref{fig:roofline-gap} expresses this rate as a fraction of each GPU's peak HBM throughput at batch size $B=128$ and context length $KV=32\mathrm{K}$.

FlashInfer BF16 exceeds $90\%$ of peak on both H100 and B200.
On B200, hardware-native TensorRT-LLM NVFP4 reaches $72.4\%$ of peak in the 4-bit group.
Among integer 4-bit baselines, BitDecoding is highest at $70.0\%$ on H100 and $36.1\%$ on B200, while its INT2 kernel reaches only $30.7\%$ and $14.5\%$, respectively.
\sys{} reaches $91.3\%$ and $60.6\%$ at INT4 and INT2 on H100, and $68.6\%$ and $44.4\%$ on B200.
Although low-bit baselines read fewer KV-cache bytes, their effective KV access throughput remains far below peak because kernel latency is bounded by the slowest pipeline stage.
\Cref{sec:analysis-issuerate} shows how operand reconstruction can demand more instruction-issue throughput than the GPU provides at peak HBM throughput.

\subsection{Reconstruction consumes instruction-issue slots}
\label{sec:analysis-issuerate}

HBM transfers, operand preparation, softmax, and Tensor Core products all require instructions from the warp schedulers.

\begin{figure}[t]
  \centering
  \subfloat[Peak throughput across GPU generations, normalized to H100 (labels show Rubin).\label{fig:machine-balance}]{%
    \includegraphics[width=0.49\columnwidth]{assets/machine-balance.pdf}}
  \hfill
  \subfloat[Instruction-issue demand per packed KV byte on H100 and B200. INT$\to$FP for converting packed integers to floating-point values; Affine is for apply scale and offset.\label{fig:ops-per-byte}]{%
    \includegraphics[width=0.49\columnwidth]{assets/ops-per-byte.pdf}}
  \caption{GPU throughput scaling and per-byte instruction-issue demand for operand construction and affine evaluation (\sys{} shown at four query heads per KV head and group size $128$).
  }
  \label{fig:reconstruction-demand}
\end{figure}


From H100 to B200, peak HBM throughput grows $2.39\times$, while peak instruction-issue throughput grows only $11\%$~\cite{nvidia2025h100spec,nvidia2025dgxb200,nvidia2025hgxb200,nvidia2026cudabpg}.
Peak dense Tensor Core throughput grows $4.55\times$ when comparing H100 FP8 with B200 FP4~\cite{nvidia2025h100spec,nvidia2025dgxb200,nvidia2025hgxb200}.
\Cref{fig:machine-balance} compares these peak throughputs across GPU generations\footnote{Rubin's figures are preliminary; we infer its instruction-issue throughput from FP32 CUDA Core throughput assuming 128 lane-slots per SM cycle~\cite{nvidia2023gpuperformance,nvidia2026rubin,nvidia2026cudabpg}.}.

Narrower stored values reduce HBM traffic without proportionally reducing reconstruction work.
We count the instructions to convert packed integers to BF16, evaluate $sq+z$ in BF16, and convert the results to FP8 or FP4 when the Tensor Core instructions require those formats.
One lane-slot corresponds to one instruction issued for one warp lane~\cite{nvidia2026cudaguide}.
Let $a_{\mathrm{issue}}$ denote the lane-slots required to construct an operand and evaluate $sq+z$ for one quantized value, $C_{\mathrm{issue}}$ the GPU's peak instruction-issue throughput in lane-slots/s, and $B_{\mathrm{HBM}}$ its peak HBM throughput in bytes/s.
Each quantized value occupies $B_{\mathrm{elem}}=0.5$~bytes in INT4 and $0.25$~bytes in INT2, excluding scales and offsets.
Streaming the cache at peak HBM throughput remains within peak instruction-issue throughput only when
\begin{equation}
  \underbrace{\frac{a_{\mathrm{issue}}}{B_{\mathrm{elem}}}}_{\text{issue demand per byte}}
  \leq
  \underbrace{\frac{C_{\mathrm{issue}}}{B_{\mathrm{HBM}}}}_{\text{peak issue per byte}}
  \label{eq:issue-balance}
\end{equation}
\Cref{fig:ops-per-byte} compares this per-byte issue demand against the hardware peak from \Cref{eq:issue-balance} on H100 and B200.

Preparing four BF16 Tensor Core operands on H100 requires two \texttt{SHF} and two \texttt{LOP3} instructions for INT$\to$FP conversion.
Two BF16 \texttt{HFMA2} instructions then evaluate the affine transformation $sq+z$, each producing two BF16 values in a 32-bit register~\cite{frantar2025marlin}.
The total is six instructions, or six lane-slots per warp lane.
Four values occupy two bytes in INT4 or one byte in INT2, giving $6/2=3$ or $6/1=6$ lane-slots/byte in~\Cref{fig:ops-per-byte}, respectively.
Therefore, reconstruction alone can consume $30\%$ and $60\%$ of H100's peak hardware budget of $9.99$ lane-slots/byte, leaving little room for other parts of the kernel.

On B200, two \texttt{F2FP} instructions convert the four BF16 values directly to FP8 or FP4 and pack them into registers.\footnote{Before PTX ISA 9.1~\cite{nvidia2026ptx}, the conversion uses FP32 intermediates and takes six instructions per four values with CUDA 13.0.
Reconstruction then costs twelve instructions, or $6$ and $12$ lane-slots/byte for INT4 and INT2, respectively.}
Reconstructing four values therefore takes $6+2=8$ lane-slots, or $8/2=4$ lane-slots/byte for INT4 and $8/1=8$ lane-slots/byte for INT2.
At peak HBM throughput, INT2 reconstruction \emph{exceeds} B200's hardware peak of $4.65$ lane-slots/byte, and INT4 reconstruction alone consumes $86\%$ of that peak.

\sys{} removes per-value affine evaluation and never materializes a reconstructed K or V tile.
Instead, it constructs Tensor Core operands directly from $q$, applies each group's scale to the other matrix input or partial product, and adds the offset as an output correction.
At four query heads per KV head and group size $128$, affine folding reduces issue demand in \Cref{fig:ops-per-byte} to $2.13$ and $4.25$ lane-slots/byte (INT4 and INT2) on H100, and to $0.88$ and $1.00$ on B200.
